\cotefile{12}
\scanpage{18}
\mspage{16}
Considérons $\Gamma=\underline{\mathrm{Aut}}_{\mathbb Q}(V_0)$, groupe algébrique % L1
sur $\mathbb Q$, et les deux hom % L2
\[
\deleted{\unsure{j_1,j_2}}i_1,i_2:\ G_{m\,\mathbb C}\longrightarrow\Gamma_{\mathbb C} % L3
\]
définissant la bigraduation de $V$ % L4
\[
\begin{aligned}&\deleted{\unsure{g_{\lambda}(x)|i_1(\lambda)}}\\[-2pt]&i(\lambda,\mu)\unsure{\longrightarrow|\mapsto|\text{est}}\lambda^p\mu^q\quad\text{sur }V^{pq}\end{aligned} % L5
\]
Soit $G$ le plus petit ss-groupe algébrique de \inserted{$\Gamma$} % L6
\deleted{tel} que $G_{\mathbb C}$ contienne $\mathrm{Im}\,i_1,\mathrm{Im}\,i_2$ ; % L7
\deleted{Alors} c'est aussi le ss-groupe algébrique de % L8
\deleted{\unsure{$G$}} $\Gamma_{\mathbb C}$ \deleted{\unsure{$G$}} engendré par les $i_1^\sigma$, $\sigma$ parcourant % L9
les \deleted{\unsure{en}} autom de $\mathbb C$. C'est un groupe % L10
algébrique connexe, appelé \uline{groupe de Mumford-Tate} % L11
de la structure de Hodge. \deleted{\unsure{Pour toute} \unsure{à quelque} % L12
représentation de $G$ \unsure{qui soit|qui} \unsure{par} sur un vectoriel % L13
$M$ de rang fini sur $\mathbb Q$, on} Si $\Sigma$ est % L14
une opération tensorielle en une variable, % L15
de sorte que $G$ opère sur $\Sigma(V_0)$, % L16
on voit que les élémts de $\Sigma(V_0)$ invariants % L17
sous $G$ sont \deleted{\unsure{exactt}} ceux qui sont % L18
invariants par \inserted{\unsure{les $i$.|les ss.|les homs}} $i_1,i_2$, i.e. qui sont de % L19
bidegré $0,0$. Cela permet de voir que % L20

% archive sheet number: 17
